\section{Proposal 3: Attribution-guided policy improvement}\label{sec:future:train}

\begin{multicols}{2}
\paragraph{Motivation} The datasets of this survey are built to score \emph{attributors}. The obvious consumer of a good attributor is the agent that produced the failure: a label $(a^\star,k^\star)$ says which parameters to move and from which state to try again. A literature search on this direction (notes in \code{notes/attribution\_as\_training\_signal.md}) gives three findings.
\begin{itemize}[leftmargin=1.1em]
\item \emph{The mechanism exists, under other names.} The strongest 2026 agent-RL methods locate a ``pivotal'', ``critical'' or ``irrecoverable'' step of a failed rollout and use it to restart rollouts there \cite{tian2026mirl,guo2026pivoarl,shi2026r3l,liu2026sprout,liang2026elpo}, to reweight the advantage at that step \cite{zhu2026fault,huang2026etro,jung2026prover,agentevolver2025,zhu2026lvrl} or to build preference pairs there \cite{li2026cso,yuan2025agentr,steca2025,cuarecovery2026}. None cites a failure-attribution benchmark as a source of signal, and none uses an attributor trained on one.
\item \emph{The benchmarks never close the loop.} Every downstream experiment inside an attribution paper is inference-time: AgenTracer feeds its tracer's reasoning back as a prompt for three rounds (gains of 4.8--14.2\%) \cite{agentracer2026}, AgentDebug re-rolls from the critical step with guidance \cite{agentdebug2026}, SearchAuditor resumes at the critical step with a repair directive \cite{searchauditor2026}, SymTrace replays one node \cite{luan2026symtrace}. No weights change. The only multi-agent blueprint---first-error localisation followed by preference pairs on the blamed agent---is a position paper without experiments \cite{yang2025blame}.
\item \emph{Attribution accuracy is never related to training value.} Nobody has asked whether a step metric from Sections~\ref{sec:g1}--\ref{sec:g4} predicts how much the agent improves when the label is used. The one direct measurement is negative: on 150 SWE-bench Verified runs, LLM-judge ``first mistake'' labels fall where the success probability does \emph{not} drop in 66\% of cases, only 7 of 188 judge answers name the decisive step, and restarting from the judge's step changes expected success by $-0.03$ to $-0.08$, whereas restarting from the certified point of no return raises it from 0.44 to 0.53 \cite{shen2026epv}. SearchAuditor's ablation adds that a generic hint at the predicted step repairs fewer failures (5.6\%) than an unguided retry (9.4\%), against 17.4\% with the auditor's directive \cite{searchauditor2026}: the location is necessary, not sufficient.
\end{itemize}
The pattern shared by the methods that work is \emph{propose, verify, then credit}: a judge or the policy proposes a step, a few current-policy continuations confirm that the outcome actually turns there, and only then does the step shape credit or data \cite{liang2026elpo,jung2026prover,li2026cso,chu2026agentgrad}. That verification step is the screen of Proposal~1 with $N$ between 1 and 8.

\paragraph{Closest prior work} Table~\ref{tab:entry} places the existing methods on the entry points defined below. Four groups:
\begin{itemize}[leftmargin=1.1em]
\item \emph{Restart methods} (MiRL, PivoARL, R$^3$L, SPROUT, ELPO) start the next rollouts from a chosen prefix; only ELPO verifies the prefix with rollouts, by binary search for the first irrecoverable step \cite{liang2026elpo}.
\item \emph{Credit methods} (FAULT, ETRO, ProVer, AgentEvolver, LVRL) reweight or truncate the advantage; ProVer verifies with $N=8$ continuations but uses positive segments only \cite{jung2026prover}.
\item \emph{Data methods} (CSO, Agent-R, CUA-Recovery) build a preference pair or a spliced correction at the step; CSO checks the alternative by policy continuation \cite{li2026cso}.
\item \emph{Multi-agent credit} (MAPPA, C3, AgentGrad) assigns per-agent credit from a coach, from counterfactual continuation of every message, or from sequential intervention, on two- to three-agent mathematics or on prompts only \cite{li2026mappa,chen2026c3,chu2026agentgrad}; nothing uses (agent, step) blame on a GAIA-style system.
\end{itemize}

\paragraph{Setting} The agent is a trainable policy $\pi_\theta$ (single- or multi-agent, each agent $a\in\mathcal{A}$ owning a parameter subset $\theta_a$) run on tasks with a success predicate $S$; the notation is that of Section~\ref{sec:future:prelim}.
\begin{itemize}[leftmargin=1.1em]
\item A group-relative outcome-only update (GRPO and its agentic variants \cite{feng2025gigpo}) applies to every step of a failed trajectory the same negative advantage $A_\tau<0$:
\[g^{\mathrm{unif}}_\tau=A_\tau\sum_{t=1}^{T}\nabla_\theta\log\pi_\theta(u_t\mid\tau_{<t}).\]
\item Under Assumption~\ref{as:absorb} only the steps in the error set are causal, so this estimator pushes down $T-|E_\alpha(\tau)|$ correct decisions per failed run; with the error rates of Table~\ref{tab:contam} that is 80--99\% of its negative gradient mass.
\item An \emph{attributor} $\mathcal{A}$ maps a trajectory to a non-negative weight vector $\hat w_\tau\in\mathbb{R}^T$, an estimate of where the causal mass lies: a point estimate $\hat w_\tau=\mathbb{1}[t=\hat k]$, a posterior $\hat w_\tau(t)=P(t^\star{=}t\mid S{=}0,\tau)$ from the hazard model of Definition~\ref{def:hazard}, or a graded profile $\hat w_\tau\propto\hat e_\tau$ from the screen of Definition~\ref{def:fa}.
\end{itemize}

\begin{definition}{Attribution-weighted update}{attrgrad}
Given $\hat w_\tau$ and a mixing weight $\lambda\in[0,1]$, the attribution-weighted policy gradient of a failed trajectory is
\[g^{\hat w}_\tau=A_\tau\sum_{t=1}^{T}\big[(1-\lambda)+\lambda\,T\hat w_\tau(t)/\|\hat w_\tau\|_1\big]\nabla_\theta\log\pi_\theta(u_t\mid\tau_{<t}),\]
so that $\lambda=0$ recovers $g^{\mathrm{unif}}_\tau$ and $\lambda=1$ with the true $w_\tau=\mathbb{1}[t=k^\star]$ concentrates the whole penalty on the decisive step. Successful trajectories are unchanged. The bias of $g^{\hat w}$ relative to the ideal $g^{w}$ is governed by the attributor's precision at the chosen tolerance---what the step metrics of this survey measure---and its variance reduction relative to $g^{\mathrm{unif}}$ by $1-\varepsilon$.
\end{definition}

\begin{definition}{Entry points}{entry}
The weight vector can enter training in four places (Fig.~\ref{fig:entry}):
\begin{itemize}[leftmargin=1.1em]
\item \textbf{R --- restart.} Rollouts for the next update start from $\tau_{<k}$ with $k\sim\rho_\tau\propto\hat w_\tau$ instead of from the task prompt. This reshapes the state distribution toward decisions that matter and lets a group of continuations share a prefix, so GRPO's within-group comparison becomes a step-level comparison at $k$. Needs Assumption~\ref{as:replay}.
\item \textbf{C --- credit.} The update of Definition~\ref{def:attrgrad}, optionally with \emph{truncation}: no dense supervision after $\hat k$, because later steps are conditioned on a state the current policy would not reach \cite{huang2026etro,guo2026pivoarl}. Needs only the log and an attributor.
\item \textbf{D --- data.} A preference pair or a spliced correction at $\hat k$: the policy's action against an alternative from which a continuation succeeds \cite{li2026cso}, or the failed prefix followed by a successful sibling branch \cite{yuan2025agentr}. Needs one rollout from $\tau_{<\hat k}\oplus\delta(s_{\hat k})$ with $\delta\in\Delta_{\mathrm{fix}}$.
\item \textbf{M --- culprit-only update} (multi-agent). $\hat w$ is indexed by (agent, step); the gradient is applied to $\theta_{a^\star}$ and the advantages of the other agents are masked, the weight-space analogue of intervention-guided prompt optimisation \cite{chu2026agentgrad}.
\end{itemize}
\end{definition}

\begin{definition}{Training value of an attributor}{trainval}
For a fixed recipe, budget $B$ of environment steps and seeds, let $\pi^{\mathcal{A}}_B$ be the policy after training with weights from $\mathcal{A}$ and $\pi^{\mathrm{unif}}_B$ the outcome-only policy at the same budget. The \emph{training value} is $V_B(\mathcal{A})=J(\pi^{\mathcal{A}}_B)-J(\pi^{\mathrm{unif}}_B)$, with $J$ the held-out success rate; the \emph{accuracy-to-value curve} is $V_B$ plotted against the attributor's step accuracy at the tolerance used by the recipe.
\end{definition}

\emph{The attributor ladder.} The experiments hold the training recipe fixed and vary only $\mathcal{A}$, ordered by cost and by how much of the outcome it has seen:
\begin{enumerate}[leftmargin=1.4em,itemsep=1pt]
\item \emph{position priors}: last decision step, per-source histogram (Proposal~2);
\item \emph{outcome variance, no attribution}: start states where the reference policy sometimes fails \cite{pivotrl2026};
\item \emph{the policy's own reflection} \cite{guo2026pivoarl};
\item \emph{prompted judge}: the all-at-once protocol of Who\&When with a frontier model;
\item \emph{trained attributors}: an 8B model trained on TracerTraj, Aegis or Who\&When~Pro with the AgenTracer recipe \cite{agentracer2026}, and the hazard model of Proposal~2 fitted at L1 on the same system's runs;
\item \emph{rollout-verified}: any of 3--5 followed by the screen of Proposal~1 with $N\le8$ at the proposed step and its neighbours \cite{liang2026elpo,jung2026prover};
\item \emph{oracle}: the rollout-measured $k^\star$ with large $N$, affordable only on a test subset.
\end{enumerate}
Each rung is scored twice: by the attribution metrics of this survey on held-out labels of the same system, and by its training value.
\end{multicols}

\begin{figure}[h]\centering
\begin{tikzpicture}[x=1cm,y=1cm,font=\scriptsize,>=Stealth]
\def\dx{0.85}
\foreach \i in {1,...,4}{\node[draw=rule,fill=surface,minimum width=0.68cm,minimum height=0.4cm,inner sep=1pt] (s\i) at (\i*\dx,0) {$s_{\i}$};}
\node[draw=critical,fill=critical!10,minimum width=0.68cm,minimum height=0.4cm,inner sep=1pt] (s5) at (5*\dx,0) {$s_{\hat k}$};
\foreach \i in {6,7}{\node[draw=rule,fill=surface,minimum width=0.68cm,minimum height=0.4cm,inner sep=1pt] (s\i) at (\i*\dx,0) {$s_{\i}$};}
\node[right=2pt of s7,text=ink2] {failed run, $S=0$};
\foreach \i/\h in {1/0.04,2/0.06,3/0.05,4/0.12,5/0.58,6/0.07,7/0.04}{\fill[s1!70] (\i*\dx-0.17,-1.0) rectangle (\i*\dx+0.17,-1.0+\h);}
\draw[rule!80!black] (0.5,-1.0) -- (7*\dx+0.4,-1.0);
\node[anchor=east,text=ink2] at (0.45,-0.75) {attribution $\hat w_\tau(t)$};
\coordinate (p) at (5*\dx,-1.08);
\node[draw=s3,fill=s3!8,rounded corners=2pt,inner sep=3pt,text width=4.6cm,align=left] (R) at (5*\dx-5.3,-2.55) {\textbf{R --- restart.} Next rollouts start from $\tau_{<k}$ with $k\sim\rho_\tau\propto\hat w_\tau$. Needs state restore.};
\node[draw=s2,fill=s2!8,rounded corners=2pt,inner sep=3pt,text width=4.6cm,align=left] (C) at (5*\dx,-2.55) {\textbf{C --- credit.} Advantage of step $t$ scaled by \mbox{$(1-\lambda)+\lambda T\hat w_\tau(t)$}, truncated after $\hat k$. Needs only the log.};
\node[draw=s7,fill=s7!8,rounded corners=2pt,inner sep=3pt,text width=4.6cm,align=left] (D) at (5*\dx+5.3,-2.55) {\textbf{D --- data.} Pair $(u_{\hat k},\delta(u_{\hat k}))$, kept if a rollout from $\delta(u_{\hat k})$ succeeds. Needs one verifying rollout.};
\draw[->,ink3] (p) -- (R.north);
\draw[->,ink3] (p) -- (C.north);
\draw[->,ink3] (p) -- (D.north);
\end{tikzpicture}
\caption{Entry points of an attribution estimate $\hat w_\tau$ into agent training (Definition~\ref{def:entry}). All three consume the same object and differ in what they need from the environment. In a multi-agent system the same vector, indexed by agent and step, also selects whose parameters are updated (entry point M).}
\label{fig:entry}
\end{figure}

\begin{table}[h]\centering\scriptsize
\setlength{\tabcolsep}{3pt}
\begin{tabularx}{\linewidth}{@{}l >{\RaggedRight\hsize=1.05\hsize}X l >{\RaggedRight\hsize=0.95\hsize}X l@{}}\toprule
Method & Attributor (who proposes $\hat k$) & Verified? & Entry & Gain vs.\ outcome-only \\\midrule
MiRL \cite{tian2026mirl} & 27B judge, clustered into 8 misbehaviours & replay only & R + judge-gated reward & TB-Lite 51.8 vs 43.4 \\
PivoARL \cite{guo2026pivoarl} & the policy's own reflection & no & R + C (prefix credit) & 83.9 vs 73.9 (GiGPO) \\
R$^3$L \cite{shi2026r3l} & the policy's own diagnosis & no & R + C (prefix mask) & ALFWorld 0.72 $\to$ 0.93 \\
SPROUT \cite{liu2026sprout} & policy as reviewer, with test diagnostics & no & R (snapshots) & TB-2.1 64.0 vs 39.3 \\
ELPO \cite{liang2026elpo} & binary-search rollouts: first irrecoverable step & MC & R + C (suffix clip) & 60.7 vs 51.2 \\
FAULT \cite{zhu2026fault} & self-diagnoser with step and quote, evidence-checked & quote match & C (priced penalties) & ALFWorld 91.0 vs 49.3 \\
ETRO \cite{huang2026etro} & 27B writer: earliest unrecovered decisive error & no & C (truncated distillation) & WebShop 85.8 vs 73.7 \\
ProVer \cite{jung2026prover} & agentic judge proposes a segment & MC, $N{=}8$ & C (positive only) & $+7$ to $+10\%$ rel. \\
AgentEvolver \cite{agentevolver2025} & judge labels every step good/bad & no & C ($\alpha{=}0.1$--$0.2$) & 47.6 vs 41.3 \\
LVRL \cite{zhu2026lvrl} & frozen LLM value function, credit = value drop & no & C & up to $+13.6$ pp \\
CSO \cite{li2026cso} & PRM candidates, expert alternative & MC, policy continuation & D (DPO at $\hat k$) & GAIA 49.5 vs 38.9 \\
Agent-R \cite{yuan2025agentr} & the policy's own first-error judgement & no & D (splice) & 70.7 vs 62.4 (splice at end) \\
CUA-Recovery \cite{cuarecovery2026} & gold per-step lineage, human-checked & yes & D (recovery SFT) & 35.2 vs 23.6 (token-matched) \\
MAPPA \cite{li2026mappa} & LLM coach scores every (agent, action) & no & M (dense) & no shared-reward ablation \\
C3 \cite{chen2026c3} & counterfactual continuation of every message & MC & M (dense) & MATH500 82.8 vs 74.5 \\
AgentGrad \cite{chu2026agentgrad} & sequential intervention: first agent whose fix flips $S$ & yes & M (prompt only) & $+11.8$ vs $+9.2$ (GEPA) \\\bottomrule
\end{tabularx}
\caption{Existing methods placed on the entry points of Definition~\ref{def:entry}. ``Verified'' states whether the proposed step is checked against the outcome before it is used. Numbers are the papers' own and are not compute-matched across rows; see \code{notes/lit\_training\_signal/} for the per-paper caveats.}
\label{tab:entry}
\end{table}

\begin{rqbox}{3}
\begin{description}[leftmargin=3.2em,style=sameline,font=\bfseries,itemsep=1.5pt]
\item[RQ\,3.1] How does the training value of an attributor depend on its attribution accuracy, and below which precision does attribution-weighted training become worse than outcome-only training at the same budget?
\item[RQ\,3.2] Which error of an attributor is the damaging one for training: a label shifted toward the step where the failure becomes visible, or a label that is randomly wrong?
\item[RQ\,3.3] Does negative credit concentrated on a rollout-verified decisive step beat unverified judge credit at equal budget, and does it avoid the reward-hacking signature of fixed process rewards?
\item[RQ\,3.4] In a multi-agent system, does updating only the blamed agent beat a shared reward, when the blame comes from an attributor trained on the datasets of this survey?
\item[RQ\,3.5] When a policy is trained against attributed failures, do the targeted failure modes shrink without the other modes growing?
\end{description}
\end{rqbox}

\begin{multicols}{2}
\paragraph{Experiments}
\begin{description}[leftmargin=1.4em,style=nextline,itemsep=1pt]
\item[\EXP{3.1} --- The accuracy-to-value curve (answers \RQ{3.1}, \RQ{3.2}).] Recipe: entry points R and C (restart from $\hat k$, prefix-only credit, truncation after $\hat k$) on a GiGPO-style learner \cite{feng2025gigpo}, 4B policy, matched rollout budget. Every rung of the ladder is trained, plus \emph{synthetically degraded} attributors: the oracle label shifted forward by $d$ steps toward the symptom, and the oracle mixed with random steps at precision $p$. The output is the curve of Definition~\ref{def:trainval}, its slope, and the precision below which attribution hurts. The evidence of \cite{shen2026epv} predicts that rung 4 sits below rung 2 on SWE-bench and that the shift $d$, not the precision $p$, is the damaging error.
\item[\EXP{3.2} --- Verified negative credit (answers \RQ{3.3}).] Entry point C alone, so that nothing is restarted. Compare (i) unverified judge weights \cite{agentevolver2025}, (ii) evidence-checked self-blame \cite{zhu2026fault}, (iii) the hazard posterior of Proposal~2 and (iv) the rollout-verified profile $\hat e_\tau$ of Proposal~1 used as $\hat w$, each at fixed $\lambda$ and at $\lambda$ annealed toward 0. ProVer uses only positive segments \cite{jung2026prover}; this experiment is its negative mirror on natural failures.
\item[\EXP{3.3} --- Culprit-only updates in a multi-agent system (answers \RQ{3.4}).] Entry point M on the MASLab systems of row 4 of Table~\ref{tab:engine} (Aegis wrappers, verifiable answers) and, for transfer, Magentic-One on GAIA (row 6, log-only). Blame from the trained attributor (rung 5) and from counterfactual message replay \cite{chen2026c3}. Conditions: shared reward to all agents, per-agent leave-one-out credit \cite{li2026sharp}, repair-every-agent data \cite{zhao2025sirius}, culprit-only gradient and culprit-only preference pairs \cite{yang2025blame}.
\item[\EXP{3.4} --- Closing the loop (answers \RQ{3.5}).] After each run of \EXP{3.1}--\EXP{3.3}, the failures of the trained policy are re-attributed and re-labelled with the failure-mode taxonomies of Section~\ref{sec:g3} (MAST modes through the Aegis labeller, AgentErrorTaxonomy). Two numbers: how much the targeted modes shrink, which MiRL reports as 43--57\% on unseen prefixes \cite{tian2026mirl}, and how much the other modes grow---the ``whack-a-mole'' risk no paper reports.
\end{description}
\end{multicols}

\medskip\noindent\begin{minipage}{\linewidth}\centering\scriptsize
\setlength{\tabcolsep}{3pt}
\begin{tabularx}{\linewidth}{@{}l >{\RaggedRight\hsize=1.25\hsize}X >{\RaggedRight\hsize=0.95\hsize}X >{\RaggedRight\hsize=0.8\hsize}X@{}}\toprule
Exp. & Data & Metrics & Baselines / reference values \\\midrule
\EXP{3.1} & \emph{Systems}: rows 1--3 of Table~\ref{tab:engine} (ALFWorld and WebShop with the AgentDebug agent, $\tau^2$-Bench, SWE-bench with a snapshotting harness); \emph{accuracy reference}: released human step labels on the same systems (AgentErrorBench 200, TrajErrBench 486, TraceElephant 44); three seeds per cell & held-out success rate $J$; training value $V_B$ (Definition~\ref{def:trainval}) at matched environment steps; the accuracy-to-value curve and its slope; the break-even precision; attribution accuracy of every rung (exact, $\pm k$) on the held-out labels & outcome-only GiGPO at the same budget; variance-selected start states \cite{pivotrl2026}; recovery training forked from the end of failed runs with no localisation \cite{reterm2026}; last-node and random-node restart \cite{luan2026symtrace} \\
\EXP{3.2} & same systems; natural failures of the current policy; weights from four attributors & $J$ and $V_B$; the \emph{reward-hacking signature}: the attributor's own score rising while $J$ falls \cite{liu2026vppo,choudhury2025agentprm,zhao2026credithurts}; fraction of negative gradient mass on the decisive step & outcome-only; unverified judge credit \cite{agentevolver2025}; ProVer-style positive-only credit \cite{jung2026prover} \\
\EXP{3.3} & \emph{Systems}: MASLab systems of row 4 (answer match, unit tests), Magentic-One on GAIA for transfer; \emph{blame reference}: Aegis-Bench (agent, mode) labels and Who\&When-style held-out labels on the same system & $J$; agent-level accuracy of the blame on held-out labels; per-agent change in generation rate; cost in rollouts per update & shared reward; leave-one-out credit \cite{li2026sharp}; repair-every-agent data \cite{zhao2025sirius}; dense per-agent coach credit \cite{li2026mappa} \\
\EXP{3.4} & failures of every trained policy of \EXP{3.1}--\EXP{3.3}, re-labelled with MAST modes (Aegis labeller) and AgentErrorTaxonomy & shrinkage of the targeted modes; growth of the other modes; net change in the mode distribution (total variation) & MiRL's 43--57\% reduction on unseen prefixes \cite{tian2026mirl}; the pre-training mode distribution \\\bottomrule
\end{tabularx}
\captionof{table}{Evaluation of Proposal 3: data, metrics and baselines per experiment. Every comparison is compute-matched on environment steps, not wall-clock, and reports both.}
\label{tab:eval_p3}
\end{minipage}\medskip

\begin{multicols}{2}
\paragraph{Evaluation: data and metrics} Table~\ref{tab:eval_p3} lists them per experiment; three conventions apply.
\begin{itemize}[leftmargin=1.1em]
\item \emph{Two scores per attributor.} Every rung is scored by the attribution metrics of this survey on held-out labels of the same system \emph{and} by its training value; the relation between the two is the result.
\item \emph{Compute matching.} Several headline gains in Table~\ref{tab:entry} come from baselines trained from the raw backbone or with fewer rollouts; here environment steps are matched and wall-clock is reported alongside.
\item \emph{Off-policy fraction.} A corrected or teacher-written step inside a prefix is off-policy for the learner; the fraction of training tokens that are off-policy is reported with every result.
\end{itemize}

\paragraph{Feasibility with the surveyed data}
\begin{itemize}[leftmargin=1.1em]
\item \emph{Environments and labels.} Rows 1--4 of Table~\ref{tab:engine} supply replayable systems with released step labels on the same system (AgentErrorBench 200, TrajErrBench 486, TraceElephant's 44 SWE-Agent traces) and Aegis-Bench's (agent, mode) labels for the multi-agent case. MiRL releases code, checkpoint and trajectories under Apache~2.0 \cite{tian2026mirl}; PivoARL and C3 release code.
\item \emph{Attributors.} AgenTracer-8B's weights are not released (\S\ref{ds:tracertraj}), so rung 5 is retrained from TracerTraj (3,998), Aegis (9,533) and Who\&When~Pro text (6,257), all ungated; the hazard model of Proposal~2 is fitted on the runs \EXP{3.1} collects anyway.
\item \emph{Budget.} One RL run at MiRL's scale is 160 updates $\times$ 8 prompts $\times$ 16 rollouts $\approx 2\times10^4$ rollouts; \EXP{3.1} with seven rungs, four degraded variants, three systems and three seeds is $\approx100$ runs, $2\times10^6$ rollouts, between $10^{10}$ and $10^{11}$ tokens with a 4B policy depending on the system (MASLab and $\tau^2$ runs are 5--20\,k tokens, SWE-bench runs up to 50\,k). The verification rollouts of rungs 6--7 add $N\le8$ continuations per failed trajectory, a small fraction of that. \EXP{3.3} on MASLab systems is cheaper (median 10 turns, 5\,k tokens).
\item \emph{Shared infrastructure.} State restore and the LLM-call cache are the components of Proposal~1's screen; the verified $(\tau,k,\delta,\hat p_\delta)$ tuples that \EXP{3.2} produces are training data for Proposal~2 at level L1$^{+}$.
\end{itemize}

\paragraph{Pitfalls to design for}
\begin{itemize}[leftmargin=1.1em]
\item \emph{Symptom shift.} Judge labels drift to the step where the failure becomes visible \cite{shen2026epv}; injected datasets place labels late (Fig.~\ref{fig:errpos}); a policy's own reflection is bounded by its own competence. Restart entry points are the most exposed, since a restart after the point of no return explores a dead state. Verify ($N\le8$) or anneal $\lambda$; never use rung 3 or 4 alone for R.
\item \emph{Reward hacking.} Any score the policy can raise without succeeding will be raised: VPPO's policy split steps to collect prefix reward \cite{liu2026vppo}, AgentPRM's process reward rose while success fell from 82\% to 70\% \cite{choudhury2025agentprm}. Guards used in the literature: conjunction with the outcome ($r=r_o\wedge r_e$ \cite{tian2026mirl}), a small or decaying weight \cite{agentevolver2025,zhao2026credithurts}, a penalty budget capped below the success--failure gap and no gradient through the diagnosis tokens \cite{zhu2026fault}.
\item \emph{Baselines that get the gain without attribution.} Variance-selected start states \cite{pivotrl2026}, recovery training forked from the \emph{end} of failed runs (Terminal-Bench 34.8 vs 28.0 \cite{reterm2026}), implicit step credit from repeated states \cite{feng2025gigpo} and SymTrace's last-node and random-node protocol \cite{luan2026symtrace}; \EXP{3.1} includes all of them as rungs or controls.
\item \emph{Off-policy prefixes.} RTI gates teacher-written steps by perplexity \cite{chen2026rti}; R$^3$L strips the guidance before training \cite{shi2026r3l}.
\item \emph{State restore.} Deterministic replay loses two thirds of terminal prefixes \cite{tian2026mirl}; per-turn snapshots lose none \cite{liu2026sprout}. Row 3 of Table~\ref{tab:engine} uses snapshots, and live-web rows cannot use entry point R at all.
\end{itemize}

\paragraph{What would count as a result}
\begin{itemize}[leftmargin=1.1em]
\item A measured accuracy-to-value curve, with the attribution precision below which attribution-weighted training is worse than outcome-only training, and the finding of whether symptom shift or random error is the damaging failure of an attributor.
\item Verified negative credit at the decisive step that beats unverified judge credit at equal budget and does not exhibit the score-up, success-down signature.
\item A culprit-only update that beats the shared reward on a multi-agent system with blame from an attributor trained on the datasets of this survey, which would be the first use of those datasets to improve an agent.
\item A closed-loop measurement showing the targeted failure modes shrink without the others growing, which is the practical meaning of ``learning from failure attribution''.
\end{itemize}
\end{multicols}
